% !TeX root = ../main.tex
\section{A Trust-Region Method with Sharper Inner-Gradient Complexity}
\label{sec:known-parameter}

We first consider the setting in which the regularity constants
$\ell$, $\mu$, and $\rho$ are available to the outer method.
The analysis has two ingredients. We begin with a fixed-scale
inexact estimate for additive universal trust-region (UTR) steps.
We then apply this estimate to the primal function using the
primal-function estimators from Section~\ref{sec:preliminaries}.
The resulting outer complexity is standard in its
$\epsilon^{-3/2}$ dependence; the main point of this section is that
the total work of the inner maximization can be accumulated without
multiplying this leading term by $\log(1/\epsilon)$.

We use
\[
    \log_+(t):=\max\{\log t,0\},
    \qquad t>0,
\]
and set $\log_+(0):=0$.


\subsection{A Fixed-Scale Inexact UTR Estimate}
\label{sec:fixed-scale-utr}

We first record the generic outer estimate used below.
Let $F:\mathbb{R}^n\to\mathbb{R}$ be twice continuously
differentiable, bounded below by
\[
    F^*:=\inf_x F(x)>-\infty,
\]
and suppose that its Hessian is $M$-Lipschitz, as in
\eqref{eq:generic-hessian-lipschitz}.

The following procedure is a certified-inexact extension of the
fixed-scale additive UTR update of \citet{jiang2026utr}. We use it
only as an outer analytical interface for the minimax method below.

\begin{assumption}[Fixed-scale inexact oracle]
\label{ass:fixed-inexact-oracle}
Fix $\epsilon>0$. At every iterate $x_k$, the oracle returns
$g_k\in\mathbb{R}^n$, a symmetric matrix
$H_k\in\mathbb{R}^{n\times n}$, and an observable certificate
$\xi_k\ge0$ satisfying
\begin{equation}
    \|g_k-\nabla F(x_k)\|
    \le
    \frac{\epsilon}{8192},
    \qquad
    \|H_k-\nabla^2F(x_k)\|
    \le
    \xi_k
    \le
    \frac{1}{64}
    \sqrt{\frac{M}{6}}\sqrt{\epsilon}.
    \label{eq:fixed-oracle-errors}
\end{equation}
\end{assumption}

Set
\begin{equation}
    h_k:=\max\{\|g_k\|,\epsilon\}.
    \label{eq:fixed-scale-h}
\end{equation}
At $x_k$, terminate if
\begin{equation}
    \|g_k\|\le\epsilon,
    \qquad
    \lambda_{\min}(H_k)
    \ge
    -\sqrt{M\epsilon}+\xi_k.
    \label{eq:fixed-scale-stopping}
\end{equation}
Otherwise, compute a global solution $d_k$ of
\begin{equation}
\begin{aligned}
    \min_{d\in\mathbb{R}^n}\quad
    &g_k^\top d
    +
    \frac12
    d^\top
    \left(
        H_k+\xi_kI
        +2\sqrt{\frac{M}{6}}\sqrt{h_k}\,I
    \right)d,
    \\
    \operatorname{s.t.}\quad
    &\|d\|
    \le
    \frac{\sqrt{h_k}}{4\sqrt{M/6}},
\end{aligned}
\label{eq:fixed-scale-subproblem}
\end{equation}
and set $x_{k+1}=x_k+d_k$.

The shift $\xi_kI$ compensates for the certified Hessian error in
both the model and the stopping test. The subproblem
\eqref{eq:fixed-scale-subproblem} retains the gradient-scaled
regularization and radius of additive UTR, while allowing inexact
first- and second-order information.

\begin{theorem}[Fixed-scale inexact UTR estimate]
\label{thm:fixed-scale-utr}
Fix $x_0$ and $\epsilon>0$, define
\[
    \Delta_F:=F(x_0)-F^*,
    \qquad
    h_0:=\max\{\|g_0\|,\epsilon\},
\]
and suppose Assumption~\ref{ass:fixed-inexact-oracle} holds.
Repeated application of
\eqref{eq:fixed-scale-stopping}--\eqref{eq:fixed-scale-subproblem}
terminates after at most
\begin{equation}
    K
    \le
    65\sqrt{\frac{M}{6}}\,
    \Delta_F\epsilon^{-3/2}
    +
    3\log_+\frac{h_0}{\epsilon}
    +2
    \label{eq:fixed-scale-complexity}
\end{equation}
nonterminal iterations, equivalently trust-region subproblem solves.
In particular,
\[
    K
    =
    O\!\left(
        \Delta_F\sqrt M\,\epsilon^{-3/2}
        +
        \log_+\frac{h_0}{\epsilon}
        +1
    \right).
\]
The returned point $\widehat x$ satisfies
\begin{equation}
    \|\nabla F(\widehat x)\|
    \le
    \frac{8193}{8192}\epsilon,
    \qquad
    \lambda_{\min}(\nabla^2F(\widehat x))
    \ge
    -\sqrt{M\epsilon}.
    \label{eq:fixed-scale-return}
\end{equation}
Consequently, replacing the internal gradient tolerance by
$8192\epsilon/8193$ gives the standard
$(\epsilon,\sqrt{M\epsilon})$ guarantee without changing the
complexity order.
\end{theorem}

The proof is given in Appendix~\ref{app:fixed-scale-proofs}.
Unlike the exact-oracle case, an inexact interior step need not
decrease $F$ monotonically. The proof controls this possible increase
through a finite-prefix argument; we defer this technical accounting
to the appendix.

For the inner-work analysis below, we need one additional consequence
of the same fixed-scale argument.

\begin{lemma}[Cubic path budget]
\label{lem:cubic-path}
Under the assumptions of Theorem~\ref{thm:fixed-scale-utr}, let
$r_k:=\|d_k\|$, and let $K$ denote the return index. Then
\begin{equation}
    M\sum_{k=0}^{K-1}r_k^3
    \le
    4\Delta_F
    +
    \frac{K}{98\,304\sqrt6}
    \frac{\epsilon^{3/2}}{\sqrt M}.
    \label{eq:cubic-path}
\end{equation}
\end{lemma}

Theorem~\ref{thm:fixed-scale-utr} controls the number of outer
subproblems. Lemma~\ref{lem:cubic-path} provides a different type of
information: it controls the cumulative motion of the outer iterates.
This path estimate is what allows the logarithmic costs of the inner
maximization to be amortized over the full trajectory rather than
bounded independently at every iteration.


\subsection{A Known-Parameter NC--SC Trust-Region Method}
\label{sec:known-parameter-method}

We now specialize the fixed-scale estimate to
\[
    P(x)=\max_y f(x,y).
\]
Recall from Lemma~\ref{lem:primal-geometry} that
\[
    M_P=4\sqrt2\,\kappa^3\rho
\]
is a valid Hessian-Lipschitz constant for $P$.
We choose the inner residual tolerance so that the corrected
estimators satisfy Assumption~\ref{ass:fixed-inexact-oracle} with
$F=P$ and $M=M_P$.

Set
\begin{equation}
    \tau_{\rm out}
    :=
    \frac{
        \mu\sqrt{M_P\epsilon/6}
    }{
        64(1+\kappa)^2\rho
    }.
    \label{eq:outer-residual-target}
\end{equation}

This target is expressed in terms of the observable residual
$\|\nabla_yf(x,y)\|$. In contrast, the MCN analysis of
\citet{luo2022finding} specifies the number of AGD iterations required
to obtain the prescribed inner accuracy.

\begin{corollary}[Residual-to-oracle certificate]
\label{cor:ncsc-fixed-oracle}
Suppose Assumption~\ref{ass:ncsc} holds and let
$M_P=4\sqrt2\,\kappa^3\rho$.
If
\begin{equation}
    \|\nabla_yf(x,y)\|
    \le
    \tau_{\rm out},
    \label{eq:outer-residual-condition}
\end{equation}
define
\begin{equation}
    \xi(x,y)
    :=
    \frac{(1+\kappa)^2\rho}{\mu}
    \|\nabla_yf(x,y)\|.
    \label{eq:ncsc-hessian-certificate}
\end{equation}
Then
\begin{equation}
    \|\widehat g(x,y)-\nabla P(x)\|
    <
    \frac{\epsilon}{8192},
    \qquad
    \|\widehat H(x,y)-\nabla^2P(x)\|
    \le
    \xi(x,y)
    \le
    \frac{1}{64}
    \sqrt{\frac{M_P}{6}}\sqrt\epsilon.
    \label{eq:ncsc-fixed-oracle-errors}
\end{equation}
Hence the fixed-scale inexact-oracle conditions
\eqref{eq:fixed-oracle-errors} hold for
$F=P$ and $M=M_P$.
\end{corollary}

The proof is deferred to Appendix~\ref{app:ncsc-oracle-calibration}.
The residual target also satisfies
\begin{subequations}
\label{eq:outer-residual-scale}
\begin{align}
    \frac{\tau_{\rm out}}{\mu}
    &=
    \Theta\!\left(
        \sqrt{\frac{\epsilon}{\kappa\rho}}
    \right)
    =
    \Theta\!\left(
        \kappa\sqrt{\frac{\epsilon}{M_P}}
    \right),
    \label{eq:outer-residual-distance-scale}
    \\
    \tau_{\rm out}
    &\ge
    \frac{\ell}{64\sqrt3}
    \sqrt{\frac{\epsilon}{M_P}}.
    \label{eq:outer-residual-lower}
\end{align}
\end{subequations}
The second relation follows directly by substituting
$M_P=4\sqrt2\,\kappa^3\rho$ and using
$4\kappa^2/(1+\kappa)^2\ge1$.

We can now state the first minimax method.

\begin{algorithm}[t]
\DontPrintSemicolon
\caption{Known-parameter NC--SC UTR with Newton-corrected derivatives}
\label{alg:ncsc-utr}
\KwIn{
$x_0,y_{-1}$, $\epsilon>0$, $\ell,\mu,\rho$, and an inner residual
module $\mathcal S$ satisfying \eqref{eq:maxres-interface}
}
$\kappa\leftarrow\ell/\mu$\;
$M_P\leftarrow4\sqrt2\,\kappa^3\rho$\;
compute $\tau_{\rm out}$ from \eqref{eq:outer-residual-target}\;
$y_0\leftarrow
\operatorname{MaxRes}_{\mathcal S}
(x_0,y_{-1},\tau_{\rm out})$\;
\For{$k=0,1,\ldots$}{
    $g_k\leftarrow\widehat g(x_k,y_k)$,
    $H_k\leftarrow\widehat H(x_k,y_k)$\;
    $\xi_k\leftarrow
    (1+\kappa)^2\rho
    \|\nabla_yf(x_k,y_k)\|/\mu$\;
    $h_k\leftarrow\max\{\|g_k\|,\epsilon\}$\;

    \If{
        $\|g_k\|\le\epsilon$ and
        $\lambda_{\min}(H_k)
        \ge-\sqrt{M_P\epsilon}+\xi_k$
    }{
        \Return{$x_k$}\;
    }

    compute a global solution $d_k$ of
    \[
    \begin{aligned}
        \min_{\|d\|\le
              \sqrt{h_k}/(4\sqrt{M_P/6})}
        \quad
        g_k^\top d
        +
        \frac12d^\top
        \left(
            H_k+\xi_k I
            +2\sqrt{\frac{M_P}{6}}\sqrt{h_k}\,I
        \right)d;
    \end{aligned}
    \]
    $x_{k+1}\leftarrow x_k+d_k$\;
    $y_{k+1}\leftarrow
    \operatorname{MaxRes}_{\mathcal S}
    (x_{k+1},y_k,\tau_{\rm out})$\;
}
\end{algorithm}

The residual used in the inner stopping condition is also used in
$\xi_k$, so no additional inner-gradient evaluation is required for
the Hessian certificate. If the inner solver happens to return a
residual smaller than $\tau_{\rm out}$, the curvature test uses the
correspondingly smaller certificate.

Algorithm~\ref{alg:ncsc-utr} does not evaluate the primal function
$P$ and does not require access to $y^*(x)$. Its outer parameters
$\tau_{\rm out}$ and $M_P$, however, depend on
$\ell$, $\mu$, and $\rho$; hence the method in this section is a
known-parameter construction even if the particular implementation of
$\mathcal S$ is itself parameter free.

\begin{theorem}[Complexity under an inner residual module]
\label{thm:ncsc-complexity}
Suppose Assumption~\ref{ass:ncsc} holds.
Fix $(x_0,y_{-1})$ and $\epsilon>0$, and define
\[
    \kappa:=\frac{\ell}{\mu},
    \qquad
    M_P:=4\sqrt2\,\kappa^3\rho,
    \qquad
    \Delta_P:=P(x_0)-P^*.
\]
Run Algorithm~\ref{alg:ncsc-utr} with any inner module satisfying
\eqref{eq:maxres-interface}. Then the algorithm returns
$\widehat x$ such that
\begin{equation}
    \|\nabla P(\widehat x)\|
    \le
    \frac{8193}{8192}\epsilon,
    \qquad
    \lambda_{\min}(\nabla^2P(\widehat x))
    \ge
    -\sqrt{M_P\epsilon}.
    \label{eq:ncsc-return}
\end{equation}

If $K$ denotes the number of nonterminal outer iterations, then
Algorithm~\ref{alg:ncsc-utr} solves exactly $K$ trust-region
subproblems and forms $K+1$ corrected-gradient and
Schur-complement Hessian estimators. Moreover,
\begin{equation}
    K
    \le
    65\sqrt{\frac{M_P}{6}}\,
    \Delta_P\epsilon^{-3/2}
    +
    3\log_+\frac{h_0}{\epsilon}
    +2,
    \label{eq:ncsc-outer-explicit}
\end{equation}
where
\[
    y_0
    :=
    \operatorname{MaxRes}_{\mathcal S}
    (x_0,y_{-1},\tau_{\rm out}),
    \qquad
    h_0
    :=
    \max\{\|\widehat g(x_0,y_0)\|,\epsilon\}.
\]
Consequently,
\begin{equation}
    K
    =
    O\!\left(
        \Delta_P
        \kappa^{3/2}\sqrt\rho\,
        \epsilon^{-3/2}
        +
        \log_+
        \frac{\|\nabla P(x_0)\|}{\epsilon}
        +1
    \right).
    \label{eq:ncsc-outer-complexity}
\end{equation}

Suppose in addition that $\mathcal S$ satisfies the accelerated
work bound \eqref{eq:inner-module-work}. Then the total number
$N_y$ of inner-gradient evaluations satisfies
\begin{align}
    N_y
    &=
    O\!\left(
        C_{\mathcal S}\sqrt\kappa\,b_{\mathcal S}
        \left[
            \Delta_P\sqrt{M_P}\,
            \epsilon^{-3/2}
            +
            \log_+
            \frac{\|\nabla P(x_0)\|}{\epsilon}
            +1
        \right]
    \right)
    \notag\\
    &\quad+
    O\!\left(
        C_{\mathcal S}\sqrt\kappa
        \log_+
        \frac{
            \|\nabla_yf(x_0,y_{-1})\|
        }{
            \tau_{\rm out}
        }
    \right).
    \label{eq:ncsc-gradient-complexity}
\end{align}
In particular, if
\[
    C_{\mathcal S}=O(1),
    \qquad
    b_{\mathcal S}=O(1+\log\kappa),
\]
then, for fixed problem data and initialization, the leading term in
$N_y$ is
\begin{equation}
    O\!\left(
        (1+\log\kappa)
        \Delta_P\kappa^2\sqrt\rho\,
        \epsilon^{-3/2}
    \right),
    \label{eq:ncsc-gradient-leading}
\end{equation}
and no $\log(1/\epsilon)$ factor multiplies this leading term.
Replacing the internal tolerance by $8192\epsilon/8193$ gives the
standard
$(\epsilon,\sqrt{M_P\epsilon})$ guarantee without changing these
orders.
\end{theorem}

The proof is given in Appendix~\ref{app:ncsc-complexity}.
The outer part follows immediately by combining
Corollary~\ref{cor:ncsc-fixed-oracle} with
Theorem~\ref{thm:fixed-scale-utr}. The nontrivial additional point is
the summation of the per-call logarithmic inner costs, which we
describe next.


\subsection{Total Inner-Gradient Complexity}
\label{sec:known-parameter-inner}

A direct worst-case estimate of every inner call would obscure the
benefit of warm starts. The proof of
Theorem~\ref{thm:ncsc-complexity} instead bounds the total logarithmic
work along the trajectory.

Suppose $y_k$ satisfies
\[
    \|\nabla_yf(x_k,y_k)\|
    \le
    \tau_{\rm out}.
\]
After the outer step $x_{k+1}=x_k+d_k$, joint
$\ell$-smoothness of $f$ gives
\begin{equation}
    \|\nabla_yf(x_{k+1},y_k)\|
    \le
    \tau_{\rm out}
    +
    \ell\|d_k\|.
    \label{eq:warm-start-residual}
\end{equation}
Hence, apart from the initial call, the logarithmic term in
\eqref{eq:inner-module-work} is bounded by
\begin{equation}
    1+
    \log_2\!\left(
        1+
        \frac{\ell\|d_k\|}{\tau_{\rm out}}
    \right).
    \label{eq:inner-log-per-call}
\end{equation}
By \eqref{eq:outer-residual-lower},
\begin{equation}
    \frac{\ell\|d_k\|}{\tau_{\rm out}}
    \le
    64\sqrt3\,
    \|d_k\|
    \sqrt{\frac{M_P}{\epsilon}}.
    \label{eq:inner-log-scale}
\end{equation}
Using
\[
    \log_2(1+t)=O(1+t^3),
    \qquad t\ge0,
\]
we obtain
\begin{equation}
    \log_2\!\left(
        1+
        \frac{\ell\|d_k\|}{\tau_{\rm out}}
    \right)
    =
    O\!\left(
        1+
        \left(
            \|d_k\|
            \sqrt{\frac{M_P}{\epsilon}}
        \right)^3
    \right).
    \label{eq:inner-log-cubic}
\end{equation}
Lemma~\ref{lem:cubic-path}, applied with $F=P$ and $M=M_P$,
therefore gives
\begin{align}
    \sum_{k=0}^{K-1}
    \left(
        \|d_k\|
        \sqrt{\frac{M_P}{\epsilon}}
    \right)^3
    &=
    \frac{M_P^{3/2}}{\epsilon^{3/2}}
    \sum_{k=0}^{K-1}\|d_k\|^3
    \notag\\
    &\le
    4\Delta_P\sqrt{M_P}\,
    \epsilon^{-3/2}
    +
    \frac{K}{98\,304\sqrt6}.
    \label{eq:inner-log-sum}
\end{align}
Thus the per-call logarithms are paid for by the same
$\epsilon^{-3/2}$-scale outer budget, rather than by multiplying the
number of outer iterations by a worst-case
$\log(1/\epsilon)$ factor.

We next verify that a standard accelerated method realizes the
abstract work condition.

For fixed $x$, consider accelerated gradient descent applied to
\[
    \phi_x(y):=-f(x,y).
\]
Starting from $y^{(0)}$, set $v^{(0)}=y^{(0)}$ and
\begin{equation}
\begin{aligned}
    y^{(j+1)}
    &=
    v^{(j)}
    +
    \frac1\ell
    \nabla_yf(x,v^{(j)}),
    \\
    v^{(j+1)}
    &=
    y^{(j+1)}
    +
    \frac{\sqrt\kappa-1}{\sqrt\kappa+1}
    \bigl(y^{(j+1)}-y^{(j)}\bigr).
\end{aligned}
\label{eq:agd-recursion}
\end{equation}

\begin{proposition}[AGD realization of the residual module]
\label{prop:agd-implementation}
Suppose Assumption~\ref{ass:ncsc} holds.
Fix $x$, an initial point $y^{(0)}$, and $\tau>0$, and define
\[
    R_{\rm in}
    :=
    \|\nabla_yf(x,y^{(0)})\|.
\]
For $\kappa>1$, the iterates in \eqref{eq:agd-recursion} satisfy
\begin{equation}
    \|y^{(N)}-y^*(x)\|
    \le
    \sqrt{\kappa+1}
    \left(
        1-\frac1{\sqrt\kappa}
    \right)^{N/2}
    \|y^{(0)}-y^*(x)\|.
    \label{eq:agd-contraction}
\end{equation}
If $R_{\rm in}\le\tau$, return $y^{(0)}$.
Otherwise, it is sufficient to take
\begin{equation}
    N
    :=
    \left\lceil
        2\sqrt\kappa
        \log_+\!\left(
            \frac{
                \sqrt{\kappa+1}\,\kappa R_{\rm in}
            }{
                \tau
            }
        \right)
    \right\rceil
    \label{eq:agd-iteration-count}
\end{equation}
and verify the residual at the terminal iterate.
Counting the initial gradient once and reusing it in the first AGD
step, this requires at most
\begin{equation}
    1+
    \left\lceil
        2\sqrt\kappa
        \log_+\!\left(
            \frac{
                \sqrt{\kappa+1}\,\kappa R_{\rm in}
            }{
                \tau
            }
        \right)
    \right\rceil
    \label{eq:agd-module-work}
\end{equation}
inner-gradient evaluations.
Consequently, AGD satisfies
\eqref{eq:inner-module-work} with
\[
    C_{\mathcal S}=O(1),
    \qquad
    b_{\mathcal S}=O(1+\log\kappa).
\]
When $\kappa=1$, one gradient step is sufficient.
\end{proposition}

The contraction in
\eqref{eq:agd-contraction} is the standard accelerated estimate for a
smooth strongly convex objective; the short verification of the
residual complexity is given in
Appendix~\ref{app:agd-realization}.

\begin{corollary}[AGD implementation of Algorithm~\ref{alg:ncsc-utr}]
\label{cor:ncsc-agd-realization}
Choose the inner module in Algorithm~\ref{alg:ncsc-utr} to be the AGD
implementation in Proposition~\ref{prop:agd-implementation}.
Then Algorithm~\ref{alg:ncsc-utr} has the outer complexity
\eqref{eq:ncsc-outer-complexity} and, in its leading term, uses
\begin{equation}
    O\!\left(
        (1+\log\kappa)
        \Delta_P\kappa^2\sqrt\rho\,
        \epsilon^{-3/2}
    \right)
    \label{eq:agd-leading-inner}
\end{equation}
inner-gradient evaluations.
No $\log(1/\epsilon)$ multiplies this leading term.
\end{corollary}

\begin{remark}[Oracle accounting]
\label{rem:known-parameter-oracle-accounting}
If Algorithm~\ref{alg:ncsc-utr} solves $K$ trust-region subproblems,
it performs $K$ nonterminal outer iterations and constructs
$K+1$ corrected gradients and Schur-complement Hessians, including
the terminal test.
The linear algebra required to apply
$(\nabla^2_{yy}f)^{-1}$ is not counted as an inner-gradient
evaluation. If these inverse actions are implemented iteratively using
Hessian-vector products, their complexity must be accounted for
separately.
The terminal evaluation of $\nabla_yf$ used to verify the residual
may be reused when forming the corrected gradient.
\end{remark}

\begin{remark}[Comparison with the raw-gradient analysis]
\label{rem:comparison-mcn}
MCN \citep{luo2022finding} uses the raw plug-in gradient
$\nabla_xf(x,y)$, for which
\begin{equation}
    \|\nabla_xf(x,y)-\nabla P(x)\|
    \le
    \ell\|y-y^*(x)\|.
    \label{eq:mcn-raw-gradient-error}
\end{equation}
Under this bound, an inner distance of order $O(\epsilon/\ell)$
is sufficient for an $O(\epsilon)$ gradient-estimation error.

More explicitly, after the summation in
\citet[Theorem~2]{luo2022finding}, the target-accuracy-dependent
logarithm in their bound involves
\begin{equation}
    \log\!\left(
        \kappa^{3/2}
        +
        \frac{\epsilon^{3/2}}
        {\rho^{3/2}\widetilde\epsilon^3}
    \right),
    \qquad
    \widetilde\epsilon
    =
    \min\left\{
        \frac{\epsilon}{192\ell},
        \frac{\sqrt{M_P\epsilon}}{48\rho}
    \right\}.
    \label{eq:mcn-explicit-log}
\end{equation}
In the small-$\epsilon$ regime, this logarithm scales as
\begin{equation}
    \log\!\left(
        \kappa^{3/2}
        +
        \left(
            \frac{\ell}{\sqrt{\rho\epsilon}}
        \right)^3
    \right).
    \label{eq:mcn-small-epsilon-log}
\end{equation}
The known quadratic accuracy of the implicit gradient
\citep{ablin2020superefficiency} instead permits an inner distance of
order $\sqrt\epsilon$.
The improvement in our total-work bound combines this accuracy with
the cumulative cubic path estimate and the summation of warm-started
inner costs. At the level of these two explicit worst-case upper
bounds, the multiplicative target-accuracy logarithm is removed
while retaining the $1+\log\kappa$ dependence of the accelerated
inner realization.
\end{remark}
